\answer
\begin{enumerate}

%%--- 章末問題 11.1 ---
\item[\ref{ex11.1}]
(a) $\tau = DT = 0.5 \times 1\,\mathrm{\mu s} = 0.5\,\mathrm{\mu s}$。
式\ref{eq11.4}より
\begin{equation*}
f_1 = \frac{1}{\pi \times 0.5\times10^{-6}} \approx 0.64\,\mathrm{MHz}，\qquad
f_2 = \frac{1}{\pi \times 5\times10^{-9}} \approx 64\,\mathrm{MHz}
\end{equation*}

(b) $f_1 < f < f_2$では$-20$\,dB，$f > f_2$では$-40$\,dBずつ下がる。

(c) $f_2 = 1/(\pi \times 50\times10^{-9}) \approx 6.4$\,MHz。
$t_r = 5$\,nsのとき$30$\,MHzはまだ$-20$\,dB/dec領域だが，
$t_r = 50$\,nsでは$6.4$\,MHz以上が$-40$\,dB/dec領域になる。
$30$\,MHzにおける包絡線は$f_2$からの追加減衰
$20\log_{10}(30/6.4) \approx 13$\,dBぶん低くなる。
立上りを緩めることが高周波ノイズの抑制に直接効く。

%%--- 章末問題 11.2 ---
\item[\ref{ex11.2}]
(a) $20\log_{10} 100 = 40$\,dB$\mathrm{\mu}$V。

(b) $10^{60/20} = 10^3\,\mathrm{\mu V} = 1$\,mV。

(c) $54 - 30 = 24$\,dBの減衰が必要。電圧比では
$10^{24/20} = 10^{1.2} \approx 16$，すなわち約$1/16$に減らすフィルタが要る。

%%--- 章末問題 11.3 ---
\item[\ref{ex11.3}]
(a) 式\ref{eq11.8}より
$L' = (4\times10^{-7}) \times 2.4 \approx 0.96\,\mathrm{\mu H/m}$，
長さ$0.2$\,mで$L \approx 0.19\,\mathrm{\mu H}$。
式\ref{eq11.9}より
$C' = \pi \times 8.85\times10^{-12} / 2.4 \approx 12\,\mathrm{pF/m}$，
長さ$0.2$\,mで$C \approx 2.3\,\mathrm{pF}$。

(b)
\begin{equation*}
f_r = \frac{1}{2\pi\sqrt{0.19\times10^{-6} \times 2.3\times10^{-12}}}
\approx 2.4\times10^{8}\,\mathrm{Hz} = 240\,\mathrm{MHz}
\end{equation*}

(c) $30$\,MHz$\sim$$1$\,GHzの\textbf{放射}エミッション試験の帯域に入る。
20\,cmの配線がこの帯域で共振し得ることは，
高周波では配線の引き回しそのものが放射ノイズの設計項目になることを意味する。

%%--- 章末問題 11.4 ---
\item[\ref{ex11.4}]
(a) Si：$dv/dt = 400/400 = 1$\,V/ns，
$i = 200\times10^{-12} \times 10^9 = 0.2$\,A。
SiC：$dv/dt = 400/40 = 10$\,V/ns，
$i = 200\times10^{-12} \times 10^{10} = 2$\,A。

(b) 電流が10倍だから$20\log_{10}10 = 20$\,dB増える。

(c) Yコンデンサは商用周波数の漏れ電流$2\pi f C V$も大地へ流すため，
感電防止の安全規格から容量の上限（数nF）が決まっており，増やせない。
そこでコモンモードだけに大きなインダクタンスとして働く
コモンモードチョークを直列に入れ，両者でフィルタを構成する（\ref{sec11.7}節）。

%%--- 章末問題 11.5 ---
\item[\ref{ex11.5}]
(a) 波形の振動周期は約$12.6$\,nsで，リンギング周波数は約$80$\,MHz。
式\ref{eq11.12}の計算値
$f_r = 1/(2\pi\sqrt{40\times10^{-9} \times 0.1\times10^{-9}}) \approx 80$\,MHz
と一致する。また抵抗が小さく減衰がごく弱いので，
ステップ応答のピークはほぼ最終値の2倍，$20$\,V近くまで達する。

(b) $t_r = 100$\,nsの立上りがつくるスペクトルの折れ点は
$f_2 = 1/(\pi \times 100\times10^{-9}) \approx 3.2$\,MHzで，
共振周波数$80$\,MHzより十分低い。励振源に$f_r$付近の成分がほとんど含まれないため，
リンギングはほぼ消える。「源を穏やかにする」対策の原理そのものである。
また減衰の時定数は$2L_p/R$で，$R = 0.1\,\Omega$では$0.8\,\mathrm{\mu s}$，
$R = 1\,\Omega$では$80$\,nsになる。振動は数周期で収まり，
リンギングのスペクトルのピークも低く・鈍くなる
（抵抗でノイズのエネルギーを熱に変えて吸収している。
フェライトコアが高周波帯で抵抗として働くのと同じ発想である）。
\end{enumerate}
